\section{Randomized --- Alg Circuit}
%%% generated by the blueprint pipeline from TCSlib.Complexity.Randomized.AlgCircuit
\subsection*{Randomized --- Alg Circuit: Introduction}

Declarations of the Lean module \texttt{TCSlib.Complexity.Randomized.AlgCircuit}.

\subsection*{Randomized --- Alg Circuit: Results}

\begin{definition}[Algebraic circuit gates]
\lean{Randomized.AlgGate}
\statementsource{Vad12}{pdf-bf2ce5422088-p016-b001}
\statementsource{BGW88}{pdf-b3f55f2a4332-p003-b004}
\statementsource{arora-barak-ch07-randomized-computation}{pdf-26d95080dbde-p004-b003}
\label{Randomized.AlgGate}
\leanok
For a fixed number $n$ of input variables, an algebraic gate is either an input variable
$x_i$ with $i<n$, the constant $1$, or the sum, difference, or product of two previously
indexed nodes, whose indices are natural numbers.
\end{definition}

\begin{definition}[Algebraic circuit]
\lean{Randomized.AlgCircuit}
\statementsource{Vad12}{pdf-bf2ce5422088-p016-b001}
\statementsource{BGW88}{
  pdf-b3f55f2a4332-p003-b004,
  pdf-b3f55f2a4332-p003-b008
}
\statementsource{arora-barak-ch07-randomized-computation}{pdf-26d95080dbde-p004-b003}
\label{Randomized.AlgCircuit}
\leanok
\uses{Randomized.AlgGate}
For a fixed number $n$ of input variables, an algebraic circuit is a finite ordered list
of algebraic gates in those variables.
\end{definition}

\begin{definition}[Evaluation of an algebraic gate]
\lean{Randomized.AlgGate.eval}
\statementsource{Vad12}{pdf-bf2ce5422088-p016-b001}
\statementsource{arora-barak-ch07-randomized-computation}{
  pdf-26d95080dbde-p004-b003,
  pdf-26d95080dbde-p005-b001
}
\label{Randomized.AlgGate.eval}
\leanok
\uses{Randomized.AlgGate}
Let $R$ be a commutative ring, let $x$ assign an element of $R$ to each of $n$ input
variables, and let $v$ be a finite list of elements of $R$. An input-variable gate
evaluates to its assigned value, the constant gate to $1$, and a gate combining nodes
$j$ and $k$ evaluates their corresponding list entries using addition, subtraction, or
multiplication, respectively; an out-of-range index contributes $0$.
\end{definition}

\begin{definition}[References to earlier circuit nodes]
\lean{Randomized.AlgGate.refsBelow}
\label{Randomized.AlgGate.refsBelow}
\leanok
\uses{Randomized.AlgGate}
For a gate with input variables or the constant $1$, the predicate is true. For an
addition, subtraction, or multiplication gate with operand indices $j$ and $k$, it is
true exactly when both indices are less than the given index $i$.
\end{definition}

\begin{definition}[Shift of gate references]
\lean{Randomized.AlgGate.shift}
\label{Randomized.AlgGate.shift}
\leanok
\uses{Randomized.AlgGate}
For a natural number offset $o$, shifting a gate leaves input variables and the constant
$1$ unchanged. In an addition, subtraction, or multiplication gate, it adds $o$ to each
operand index.
\end{definition}

\begin{theorem}[Evaluation commutes with ring homomorphisms]
\lean{Randomized.AlgGate.map_eval}
\label{Randomized.AlgGate.map_eval}
\leanok
\uses{Randomized.AlgGate, Randomized.AlgGate.eval}
\difficulty{4}
Let $R$ and $S$ be commutative rings, let $f:R\to S$ be a ring homomorphism, let $x$
assign elements of $R$ to the $n$ inputs, let $v$ be a list of elements of $R$, and let
$g$ be an algebraic gate. Applying $f$ to the value of $g$ at $(x,v)$ gives the value of
$g$ at the input assignment obtained by applying $f$ to each value of $x$, and the list
obtained by applying $f$ entrywise to $v$.
\end{theorem}

\begin{theorem}[Evaluation under shifted gate references]
\lean{Randomized.AlgGate.eval_shift}
\label{Randomized.AlgGate.eval_shift}
\leanok
\uses{Randomized.AlgGate, Randomized.AlgGate.eval, Randomized.AlgGate.shift}
\difficulty{4}
Let $x$ assign a value in a commutative ring to each of $n$ inputs, let $v$ and $w$ be
lists of ring elements, and let $g$ be an algebraic gate. If every node reference in $g$
is increased by the length of $v$, then evaluating the shifted gate on the concatenated
list $v$ followed by $w$ gives the same value as evaluating $g$ on $w$.
\end{theorem}

\begin{definition}[Execution of an algebraic circuit]
\lean{Randomized.AlgCircuit.run}
\statementsource{arora-barak-ch07-randomized-computation}{
  pdf-26d95080dbde-p004-b003,
  pdf-26d95080dbde-p005-b001
}
\label{Randomized.AlgCircuit.run}
\leanok
\uses{Randomized.AlgGate, Randomized.AlgGate.eval}
Let $R$ be a commutative ring, let $x$ assign an element of $R$ to each of $n$ input
variables, and let $g_1,\ldots,g_m$ be an ordered list of algebraic gates. Running the
list produces, in order, the values of each gate, with each gate evaluated using the
input assignment and the values of all earlier gates.
\end{definition}

\begin{definition}[Algebraic circuit size]
\lean{Randomized.AlgCircuit.size}
\statementsource{arora-barak-ch07-randomized-computation}{pdf-26d95080dbde-p005-b003}
\label{Randomized.AlgCircuit.size}
\leanok
\uses{Randomized.AlgCircuit}
The size of an algebraic circuit is the number of gates in its ordered list.
\end{definition}

\begin{definition}[Algebraic circuit evaluation]
\lean{Randomized.AlgCircuit.eval}
\statementsource{Vad12}{pdf-bf2ce5422088-p016-b001}
\statementsource{arora-barak-ch07-randomized-computation}{
  pdf-26d95080dbde-p004-b003,
  pdf-26d95080dbde-p005-b001
}
\label{Randomized.AlgCircuit.eval}
\leanok
\uses{Randomized.AlgCircuit, Randomized.AlgCircuit.run, Randomized.AlgCircuit.size}
Let $R$ be a commutative ring, and let $C$ be an algebraic circuit in $n$ variables with
input assignment $x\in R^n$. Its value is the value of the last gate in the list
produced by running the circuit; if the circuit has no gates, its value is $0$.
\end{definition}

\begin{definition}[Polynomial associated to a circuit]
\lean{Randomized.AlgCircuit.toMvPolynomial}
\statementsource{Vad12}{
  pdf-bf2ce5422088-p014-b001,
  pdf-bf2ce5422088-p016-b001
}
\statementsource{arora-barak-ch07-randomized-computation}{pdf-26d95080dbde-p004-b003}
\label{Randomized.AlgCircuit.toMvPolynomial}
\leanok
\uses{Randomized.AlgCircuit, Randomized.AlgCircuit.eval}
Every algebraic circuit in $n$ variables determines a polynomial with integer
coefficients in those variables, obtained by evaluating the circuit with each input
variable represented by its corresponding indeterminate.
\end{definition}

\begin{definition}[Well-formed algebraic circuit]
\lean{Randomized.AlgCircuit.WellFormed}
\statementsource{arora-barak-ch07-randomized-computation}{pdf-26d95080dbde-p004-b003}
\label{Randomized.AlgCircuit.WellFormed}
\leanok
\uses{Randomized.AlgCircuit, Randomized.AlgGate.refsBelow}
An algebraic circuit on $n$ inputs is well formed if its gate list is nonempty and, for
every index $i$ of a gate in the list, all node references used by that gate have
indices strictly less than $i$.
\end{definition}

\begin{theorem}[Execution of an Appended Gate]
\lean{Randomized.AlgCircuit.run_append_singleton}
\label{Randomized.AlgCircuit.run_append_singleton}
\leanok
\uses{Randomized.AlgCircuit.run, Randomized.AlgGate, Randomized.AlgGate.eval}
\difficulty{1}
Let $R$ be a commutative ring, let $x$ assign an element of $R$ to each of $n$ input
variables, let $gs$ be a finite list of arithmetic gates, and let $g$ be an arithmetic
gate. Running the list formed by appending $g$ to $gs$ gives the values from running
$gs$, followed by the value of $g$ evaluated using those values.
\end{theorem}

\begin{theorem}[Run length]
\lean{Randomized.AlgCircuit.length_run}
\label{Randomized.AlgCircuit.length_run}
\leanok
\uses{Randomized.AlgCircuit.run, Randomized.AlgCircuit.run_append_singleton, Randomized.AlgGate}
\difficulty{4}
For any input assignment and any finite list of arithmetic gates, executing the gates in
order produces a list of values with exactly as many entries as there are gates.
\end{theorem}

\begin{theorem}[Indexed lookup after appending one value]
\lean{Randomized.AlgCircuit.getD_append_singleton}
\label{Randomized.AlgCircuit.getD_append_singleton}
\leanok
\difficulty{4}
Let $R$ be a commutative ring, let $v$ be an element of $R$, let $vs$ be a finite list
of elements of $R$, and let $j$ be a natural number. Looking up index $j$ in the list
obtained by appending $v$ to $vs$, using $0$ when the index is out of range, yields
either the corresponding lookup in $vs$ or $v$.
\end{theorem}

\begin{theorem}[Evaluation under a ring homomorphism]
\lean{Randomized.AlgCircuit.run_map}
\label{Randomized.AlgCircuit.run_map}
\leanok
\uses{Randomized.AlgCircuit.run, Randomized.AlgCircuit.run_append_singleton, Randomized.AlgGate, Randomized.AlgGate.map_eval}
\difficulty{4}
Let $R$ and $S$ be commutative rings, let $f:R\to S$ be a ring homomorphism, let $x$
assign an element of $R$ to each of $n$ inputs, and let $g_1,\ldots,g_m$ be an ordered
list of arithmetic gates. Evaluating these gates successively at the inputs obtained by
applying $f$ gives the list obtained by applying $f$ to each value in the successive
evaluation of the original gates at $x$.
\end{theorem}

\begin{theorem}[Evaluation and ring homomorphisms]
\lean{Randomized.AlgCircuit.map_eval}
\label{Randomized.AlgCircuit.map_eval}
\leanok
\uses{Randomized.AlgCircuit, Randomized.AlgCircuit.eval, Randomized.AlgCircuit.run, Randomized.AlgCircuit.run_map, Randomized.AlgCircuit.size}
\difficulty{4}
Let $R$ and $S$ be commutative rings, let $f:R\to S$ be a ring homomorphism, and let $C$
be a finite arithmetic circuit on $n$ inputs. For every input $x\in R^n$, applying $f$
to the circuit's value at $x$ gives the value of the same circuit at the input obtained
by applying $f$ coordinatewise: $f(C(x))=C(f\circ x)$.
\end{theorem}

\begin{theorem}[Polynomial representation of circuit evaluation]
\lean{Randomized.AlgCircuit.aeval_toMvPolynomial}
\label{Randomized.AlgCircuit.aeval_toMvPolynomial}
\leanok
\uses{Randomized.AlgCircuit, Randomized.AlgCircuit.eval, Randomized.AlgCircuit.map_eval, Randomized.AlgCircuit.toMvPolynomial}
\difficulty{4}
Let $C$ be an arithmetic circuit with $n$ inputs, let $R$ be a commutative ring, and let
$x\in R^n$. Evaluating the integer-coefficient multivariate polynomial represented by
$C$ at $x$ gives the same value as evaluating the circuit $C$ at $x$.
\end{theorem}

\begin{theorem}[Circuit evaluation as a polynomial]
\lean{Randomized.AlgCircuit.eval_toMvPolynomial}
\label{Randomized.AlgCircuit.eval_toMvPolynomial}
\leanok
\uses{Randomized.AlgCircuit, Randomized.AlgCircuit.aeval_toMvPolynomial, Randomized.AlgCircuit.eval, Randomized.AlgCircuit.toMvPolynomial}
\difficulty{3}
Let $C$ be an algebraic circuit on $n$ inputs, and let $x$ assign an integer to each
input. Evaluating the multivariable polynomial associated with $C$ at $x$ equals the
circuit's evaluation at $x$.
\end{theorem}

\begin{theorem}[Polynomial vanishing criterion]
\lean{Randomized.AlgCircuit.toMvPolynomial_eq_zero_iff}
\label{Randomized.AlgCircuit.toMvPolynomial_eq_zero_iff}
\leanok
\uses{Randomized.AlgCircuit, Randomized.AlgCircuit.eval, Randomized.AlgCircuit.eval_toMvPolynomial, Randomized.AlgCircuit.toMvPolynomial}
\difficulty{3}
For any arithmetic circuit $C$ on $n$ inputs, its integer-coefficient multivariate
polynomial is zero if and only if its evaluation is zero at every assignment of integers
to those inputs.
\end{theorem}

\begin{theorem}[Degree bound for circuit execution]
\lean{Randomized.AlgCircuit.totalDegree_run_le}
\label{Randomized.AlgCircuit.totalDegree_run_le}
\leanok
\uses{Randomized.AlgCircuit.getD_append_singleton, Randomized.AlgCircuit.run, Randomized.AlgCircuit.run_append_singleton, Randomized.AlgGate, Randomized.AlgGate.eval}
\difficulty{6}
Let $n$ be a nonnegative integer, let $g_1,\ldots,g_m$ be a list of algebraic gates on
$n$ inputs, and evaluate them in order at the formal input variables over $\mathbb{Z}$.
For every position $j$, the total degree of the value at that position, taking the value
to be zero when the position is out of range, is at most $2^m$.
\end{theorem}

\begin{theorem}[Degree bound for circuit polynomials]
\lean{Randomized.AlgCircuit.totalDegree_toMvPolynomial_le}
\label{Randomized.AlgCircuit.totalDegree_toMvPolynomial_le}
\leanok
\uses{Randomized.AlgCircuit, Randomized.AlgCircuit.size, Randomized.AlgCircuit.toMvPolynomial, Randomized.AlgCircuit.totalDegree_run_le}
\difficulty{1}
Let $C$ be an algebraic circuit on $n$ inputs, with size equal to its number of gates.
The total degree of the polynomial computed by $C$ over $\mathbb{Z}$ is at most
$2^{\mathrm{size}(C)}$.
\end{theorem}

\begin{theorem}[Magnitude bound for circuit outputs]
\lean{Randomized.AlgCircuit.abs_run_le}
\label{Randomized.AlgCircuit.abs_run_le}
\leanok
\uses{Randomized.AlgCircuit.getD_append_singleton, Randomized.AlgCircuit.run, Randomized.AlgCircuit.run_append_singleton, Randomized.AlgGate, Randomized.AlgGate.eval}
\difficulty{5}
Let $n$ and $X$ be natural numbers with $2\le X$, let $x\in\mathbb{Z}^n$ satisfy
$|x_i|\le X$ for every input coordinate $i$, and let $G$ be any list of arithmetic gates
on these inputs. For every node index $j$, the value at that position in the circuit run
(taken to be zero when the position is absent) has magnitude at most $X^{2^{|G|}}$,
where $|G|$ is the number of gates.
\end{theorem}

\begin{theorem}[Sharper circuit evaluation magnitude bound]
\lean{Randomized.AlgCircuit.abs_eval_le_pow}
\label{Randomized.AlgCircuit.abs_eval_le_pow}
\leanok
\uses{Randomized.AlgCircuit, Randomized.AlgCircuit.abs_run_le, Randomized.AlgCircuit.eval, Randomized.AlgCircuit.size}
\difficulty{1}
Let $C$ be an algebraic circuit on $n$ integer inputs, and let $X\ge2$ be an integer. If
every input $x_i$ satisfies $|x_i|\le X$, then the absolute value of the circuit's
output is at most $X^{2^{\operatorname{size}(C)}}$.
\end{theorem}

\begin{theorem}[Circuit evaluation magnitude bound]
\lean{Randomized.AlgCircuit.abs_eval_le}
\label{Randomized.AlgCircuit.abs_eval_le}
\leanok
\uses{Randomized.AlgCircuit, Randomized.AlgCircuit.abs_eval_le_pow, Randomized.AlgCircuit.eval, Randomized.AlgCircuit.size}
\difficulty{3}
Let $C$ be an algebraic circuit on $n$ integer inputs, and let $X\ge1$ be an integer. If
every input $x_i$ satisfies $|x_i|\le X$, then the absolute value of the circuit's
output is at most $(2X)^{2^{\operatorname{size}(C)}}$.
\end{theorem}

\begin{theorem}[Execution after shifting gate references]
\lean{Randomized.AlgCircuit.run_append_shift}
\label{Randomized.AlgCircuit.run_append_shift}
\leanok
\uses{Randomized.AlgCircuit.length_run, Randomized.AlgCircuit.run, Randomized.AlgGate, Randomized.AlgGate.eval, Randomized.AlgGate.eval_shift, Randomized.AlgGate.shift}
\difficulty{4}
Let $x$ be an assignment to the inputs, and let $g$ and $h$ be finite lists of
arithmetic gates. If every gate reference in $h$ is shifted forward by the length of
$g$, then executing the concatenated list yields the values from executing $g$, followed
by the values from executing $h$: $\operatorname{run}(x,g\mathbin{+\!\!+}
\operatorname{shift}_{|g|}(h))=\operatorname{run}(x,g)\mathbin{+\!\!+}\operatorname{run}(x,h)$.
\end{theorem}

\begin{definition}[Circuit subtraction construction]
\lean{Randomized.AlgCircuit.subCircuit}
\statementsource{arora-barak-ch07-randomized-computation}{pdf-26d95080dbde-p004-b003}
\label{Randomized.AlgCircuit.subCircuit}
\leanok
\uses{Randomized.AlgCircuit, Randomized.AlgCircuit.size, Randomized.AlgGate.shift}
Given two algebraic circuits $C$ and $C'$ on the same number of inputs, form a circuit
by listing the gates of $C$, then the gates of $C'$ with their node references shifted
by the size of $C$, and then a subtraction gate whose inputs are the last nodes of these
two circuits, at indices $|C|-1$ and $|C|+|C'|-1$, respectively.
\end{definition}

\begin{theorem}[Size of a subcircuit construction]
\lean{Randomized.AlgCircuit.size_subCircuit}
\label{Randomized.AlgCircuit.size_subCircuit}
\leanok
\uses{Randomized.AlgCircuit, Randomized.AlgCircuit.size, Randomized.AlgCircuit.subCircuit}
\difficulty{3}
For any two arithmetic circuits $C$ and $C'$ on the same number of inputs, the circuit
formed by placing their gate lists together, shifting the node references of the second
list by the size of the first, and appending a subtraction gate has size
$\mathrm{size}(C)+\mathrm{size}(C')+1$.
\end{theorem}

\begin{theorem}[Evaluation of a circuit difference]
\lean{Randomized.AlgCircuit.eval_subCircuit}
\label{Randomized.AlgCircuit.eval_subCircuit}
\leanok
\uses{Randomized.AlgCircuit, Randomized.AlgCircuit.eval, Randomized.AlgCircuit.length_run, Randomized.AlgCircuit.run, Randomized.AlgCircuit.run_append_shift, Randomized.AlgCircuit.run_append_singleton, Randomized.AlgCircuit.size, Randomized.AlgCircuit.subCircuit, Randomized.AlgGate.eval}
\difficulty{4}
Let $C$ and $C'$ be algebraic circuits on the same finite set of inputs, each with at
least one gate, and let the inputs take values in a commutative ring. The circuit formed
by adjoining the shifted gates of $C'$ to those of $C$ and then subtracting their output
nodes evaluates to the value of $C$ minus the value of $C'$.
\end{theorem}

\begin{theorem}[Well-formedness of a combined circuit]
\lean{Randomized.AlgCircuit.WellFormed.subCircuit}
\label{Randomized.AlgCircuit.WellFormed.subCircuit}
\leanok
\uses{Randomized.AlgCircuit, Randomized.AlgCircuit.WellFormed, Randomized.AlgCircuit.size, Randomized.AlgCircuit.subCircuit, Randomized.AlgGate.refsBelow, Randomized.AlgGate.shift}
\difficulty{5}
Let $C$ and $C'$ be well-formed algebraic circuits on the same $n$ input variables: each
circuit has at least one gate, and every gate referring to nodes refers only to earlier
nodes. Form a circuit by listing the gates of $C$, then the gates of $C'$ with every
node reference shifted by the size of $C$, and finally a subtraction gate whose inputs
are the last nodes of the two circuits. The resulting circuit is well formed.
\end{theorem}

\begin{theorem}[Polynomial of a subcircuit construction]
\lean{Randomized.AlgCircuit.toMvPolynomial_subCircuit}
\label{Randomized.AlgCircuit.toMvPolynomial_subCircuit}
\leanok
\uses{Randomized.AlgCircuit, Randomized.AlgCircuit.eval_subCircuit, Randomized.AlgCircuit.subCircuit, Randomized.AlgCircuit.toMvPolynomial}
\difficulty{1}
Let $C$ and $C'$ be arithmetic circuits on the same finite set of inputs, each with a
nonempty gate list. The polynomial represented by the circuit formed from them using the
specified shifted-node subtraction construction is the polynomial represented by $C$
minus the polynomial represented by $C'$.
\end{theorem}

\begin{theorem}[Polynomial equality and circuit difference]
\lean{Randomized.AlgCircuit.toMvPolynomial_eq_iff_subCircuit}
\label{Randomized.AlgCircuit.toMvPolynomial_eq_iff_subCircuit}
\leanok
\uses{Randomized.AlgCircuit, Randomized.AlgCircuit.subCircuit, Randomized.AlgCircuit.toMvPolynomial, Randomized.AlgCircuit.toMvPolynomial_subCircuit}
\difficulty{3}
Let $C$ and $C'$ be arithmetic circuits on the same finite set of inputs, each with a
nonempty gate list. The integer-coefficient multivariate polynomials represented by $C$
and $C'$ are equal if and only if the polynomial represented by the circuit formed from
them using the specified shifted-node subtraction construction is zero.
\end{theorem}

\begin{definition}[Encoding arithmetic gates]
\lean{Randomized.AlgGate.encode}
\label{Randomized.AlgGate.encode}
\leanok
\uses{BoolCircuit.encodeNat, Randomized.AlgGate}
Encode each arithmetic gate as a Boolean list using a three-bit operation tag followed
by its parameters: an input variable uses tag $000$ and the unary natural-number
encoding of its index; the constant $1$ uses tag $001$; addition uses tag $010$ followed
by the encodings of its two operand indices; subtraction uses tag $011$ followed by
those encodings; and multiplication uses tag $100$ followed by those encodings. The
unary encoding of a natural number $k$ consists of $k$ true bits followed by one false
bit.
\end{definition}

\begin{definition}[Parsing a binary gate]
\lean{Randomized.AlgGate.parseBin}
\label{Randomized.AlgGate.parseBin}
\leanok
\uses{BoolCircuit.decodeNat, Randomized.AlgGate}
Given a function that constructs a gate from two natural-number indices and a Boolean
list, attempt to decode two natural numbers successively from the list. If both
decodings succeed, return the constructed gate together with the unused suffix; if
either decoding fails, return no result.
\end{definition}

\begin{definition}[Parsing an arithmetic gate]
\lean{Randomized.AlgGate.parse}
\label{Randomized.AlgGate.parse}
\leanok
\uses{BoolCircuit.decodeNat, Randomized.AlgGate, Randomized.AlgGate.parseBin}
Given a natural number $n$ and a Boolean list, parse its initial three bits as an
operation tag. Tags $000$, $001$, $010$, $011$, and $100$ specify, respectively, an
input variable, the constant $1$, addition, subtraction, and multiplication; all other
tags and incomplete tags fail. For an input-variable tag, decode one natural number and
return that variable only if its index is less than $n$; for a binary-operation tag,
decode two successive natural numbers and return the corresponding gate, in each case
preserving the unconsumed suffix. The constant tag returns the constant and the suffix
following its tag.
\end{definition}

\begin{theorem}[Parsing an Encoded Algebraic Gate]
\lean{Randomized.AlgGate.parse_encode}
\label{Randomized.AlgGate.parse_encode}
\leanok
\uses{BoolCircuit.decodeNat_encodeNat, Randomized.AlgGate, Randomized.AlgGate.encode, Randomized.AlgGate.parse, Randomized.AlgGate.parseBin}
\difficulty{4}
Fix $n\in\mathbb{N}$. Encode a variable $x_i$ by the tag $000$ followed by the unary
encoding of $i$, the constant $1$ by $001$, and addition, subtraction, and
multiplication by the tags $010$, $011$, and $100$, respectively, each followed by the
unary encodings of its two node indices. Let the parser read these tags and encodings,
accepting a variable only when its index is less than $n$; then for every algebraic gate
$g$ on $n$ inputs and every Boolean string $r$, parsing the encoding of $g$ followed by
$r$ returns $g$ and leaves $r$ unconsumed.
\end{theorem}

\begin{definition}[Encoding a list of arithmetic gates]
\lean{Randomized.AlgCircuit.encodeGates}
\label{Randomized.AlgCircuit.encodeGates}
\leanok
\uses{Randomized.AlgGate, Randomized.AlgGate.encode}
Encode a list of arithmetic gates over $n$ variables recursively: the empty list is
encoded as the single false bit, and a nonempty list is encoded as a true bit followed
by the encoding of its first gate and then the encoding of the remaining list.
\end{definition}

\begin{definition}[Gate-list parsing]
\lean{Randomized.AlgCircuit.parseGates}
\label{Randomized.AlgCircuit.parseGates}
\leanok
\uses{Randomized.AlgGate, Randomized.AlgGate.parse}
For a fixed number $n$ of input variables and a natural-number parsing bound $m$, the
gate-list parser returns either failure or a list of gates. It accepts the single-bit
list $[\mathrm{false}]$ as the empty list; within the bound, a leading $\mathrm{true}$
signals a gate to parse, followed recursively by the remaining gates, while any other
input pattern fails.
\end{definition}

\begin{definition}[Circuit encoding]
\lean{Randomized.AlgCircuit.encode}
\label{Randomized.AlgCircuit.encode}
\leanok
\uses{BoolCircuit.encodeNat, Randomized.AlgCircuit, Randomized.AlgCircuit.encodeGates}
The encoding of a circuit with $n$ input variables is the concatenation of the
natural-number encoding of $n$ and the encoding of its ordered gate list.
\end{definition}

\begin{definition}[Circuit decoding]
\lean{Randomized.AlgCircuit.decode}
\label{Randomized.AlgCircuit.decode}
\leanok
\uses{BoolCircuit.decodeNat, Randomized.AlgCircuit, Randomized.AlgCircuit.encode, Randomized.AlgCircuit.parseGates}
Given a finite Boolean list, first decode its initial natural number $n$, then parse the
remaining list as a gate list using its length as the parsing bound. The decoder returns
the resulting circuit with $n$ input variables only when re-encoding that circuit
reproduces the original list; otherwise it fails.
\end{definition}

\begin{theorem}[Length of an encoded gate list]
\lean{Randomized.AlgCircuit.length_encodeGates}
\label{Randomized.AlgCircuit.length_encodeGates}
\leanok
\uses{Randomized.AlgCircuit.encodeGates, Randomized.AlgGate}
\difficulty{4}
For any finite list $gs$ of algebraic gates on $n$ inputs, the number of gates in $gs$
is strictly less than the length of its Boolean encoding.
\end{theorem}

\begin{theorem}[Parsing an encoded gate list]
\lean{Randomized.AlgCircuit.parseGates_encodeGates}
\label{Randomized.AlgCircuit.parseGates_encodeGates}
\leanok
\uses{Randomized.AlgCircuit.encodeGates, Randomized.AlgCircuit.parseGates, Randomized.AlgGate, Randomized.AlgGate.parse_encode}
\difficulty{4}
Let $g_1,\ldots,g_m$ be a finite list of arithmetic gates, and let $f$ be a
natural-number parsing budget with $m\le f$. Parsing the Boolean encoding of this list
with budget $f$ returns exactly the original list.
\end{theorem}

\begin{theorem}[Decoding an encoded circuit]
\lean{Randomized.AlgCircuit.decode_encode}
\label{Randomized.AlgCircuit.decode_encode}
\leanok
\uses{BoolCircuit.decodeNat_encodeNat, Randomized.AlgCircuit, Randomized.AlgCircuit.decode, Randomized.AlgCircuit.encode, Randomized.AlgCircuit.encodeGates, Randomized.AlgCircuit.length_encodeGates, Randomized.AlgCircuit.parseGates, Randomized.AlgCircuit.parseGates_encodeGates}
\difficulty{4}
For every finite input arity $n$ and every arithmetic circuit $C$ with $n$ inputs,
decoding the Boolean-list encoding of $C$ returns $C$ exactly.
\end{theorem}

\begin{theorem}[Recovery of an encoded circuit]
\lean{Randomized.AlgCircuit.encode_of_decode}
\label{Randomized.AlgCircuit.encode_of_decode}
\leanok
\uses{Randomized.AlgCircuit, Randomized.AlgCircuit.decode, Randomized.AlgCircuit.encode}
\difficulty{4}
For a natural number of inputs and an algebraic circuit on those inputs, if decoding a
Boolean list returns that circuit, then encoding the circuit recovers the list.
\end{theorem}

\begin{theorem}[Characterization of circuit decoding]
\lean{Randomized.AlgCircuit.decode_eq_some_iff}
\label{Randomized.AlgCircuit.decode_eq_some_iff}
\leanok
\uses{Randomized.AlgCircuit, Randomized.AlgCircuit.decode, Randomized.AlgCircuit.decode_encode, Randomized.AlgCircuit.encode, Randomized.AlgCircuit.encode_of_decode}
\difficulty{1}
For every Boolean list $x$ and every finite-input arithmetic circuit $C$, decoding $x$
returns $C$ exactly if and only if $x$ is the Boolean-list encoding of $C$.
\end{theorem}

\begin{theorem}[Circuit encoding length]
\lean{Randomized.AlgCircuit.lt_length_encode}
\label{Randomized.AlgCircuit.lt_length_encode}
\leanok
\uses{BoolCircuit.encodeNat, Randomized.AlgCircuit, Randomized.AlgCircuit.encode, Randomized.AlgCircuit.length_encodeGates, Randomized.AlgCircuit.size}
\difficulty{3}
Let a circuit have $n$ input variables and $s$ gates. The length of its Boolean encoding
is strictly greater than $n+s$.
\end{theorem}

\begin{definition}[Zero polynomial circuit language]
\lean{Randomized.ZEROP}
\statementsource{Vad12}{pdf-bf2ce5422088-p016-b002}
\statementsource{arora-barak-ch07-randomized-computation}{pdf-26d95080dbde-p004-b003}
\label{Randomized.ZEROP}
\leanok
\uses{Randomized.AlgCircuit, Randomized.AlgCircuit.WellFormed, Randomized.AlgCircuit.encode, Randomized.AlgCircuit.toMvPolynomial}
The language consists of encodings of pairs $(n,C)$ in which $C$ is a well-formed
algebraic circuit on $n$ inputs and the polynomial computed by $C$ is zero.
\end{definition}

\begin{theorem}[Decoding Zero Polynomial Circuits]
\lean{Randomized.mem_ZEROP_iff}
\label{Randomized.mem_ZEROP_iff}
\leanok
\uses{Randomized.AlgCircuit, Randomized.AlgCircuit.WellFormed, Randomized.AlgCircuit.decode, Randomized.AlgCircuit.decode_eq_some_iff, Randomized.AlgCircuit.toMvPolynomial, Randomized.ZEROP}
\difficulty{3}
For every bit string $x$, membership in the language of encodings of nonempty,
well-formed arithmetic circuits whose associated integer-coefficient polynomial is zero
is equivalent to the existence of a circuit code that decodes from $x$, is well formed,
and has associated polynomial equal to zero.
\end{theorem}
